\documentclass[../BimodalReference.tex]{subfiles}
\begin{document}

\section{Task Semantics}

\subsection{Task Frames}

Task frames are the fundamental semantic structures for \textbf{TM}.
They abstract from universal laws governing transitions between world states while still retaining the temporal duration for a transition to complete.

The following primitives are required to define a task frame:

\begin{center}
\begin{tabular}{@{}lll@{}}
\toprule
Primitive & Type & Description \\
\midrule
$\worldstate$ & Type & World states (nonempty) \\
$D$ & Type & Temporal durations (nontrivial totally ordered abelian group) \\
$\poscone$ & $\{x : D \mid 0 \leq x\}$ & Positive cone of durations \\
$w \taskto{x} u$ & $\worldstate \to \poscone \to \worldstate \to \text{Prop}$ & Primitive task relation \\
$w \taskto{x} u$ & $\worldstate \to D \to \worldstate \to \text{Prop}$ & Extended task relation \\
$\Fib(w, x)$ & $\worldstate \to D \to \mathcal{P}(\worldstate)$ & Fiber of $w$ at duration $x$ \\
$\taskcone{w}{x}$ & $\worldstate \to \{x : D \mid 0 < x\} \to \mathcal{P}(\worldstate)$ & Cone of radius $x$ around $w$ \\
$\taskseg{w}{v}{x}{y}$ & $\worldstate \to \worldstate \to \poscone \to \poscone \to \mathcal{P}(\worldstate)$ & Segment between $w$ and $v$ \\
\bottomrule
\end{tabular}
\end{center}

\noindent
The task relation is primitive only on the positive cone; the extended relation is obtained from it by the converse convention stated below, and the fibers, cones, and segments are all defined in terms of the extended relation.

% Restates the JPL paper's \label{def:task-relation}. Verbatim source text (quoted so a
% renamed anchor stays detectable by text search):
%   "A \textit{task relation} on a nonempty set of \textit{world states} $W$ over a temporal
%    order $\D$ is any parameterized relation $w \Rightarrow_x u$ for $w,u \in W$ and
%    $x \in D^+$, extended to negative durations by the \textit{converse convention}
%    $w \Rightarrow_{-x} u \coloneq u \Rightarrow_{x} w$ for $x \geq 0$ ..."
%   "\item[\it Fiber:] $\Fib(w, x) \coloneq \set{u \in W : w \Rightarrow_x u}$."
%   "\item[\it Cone:] $(w)_x \coloneq \bigcup\limits_{\vert{y} < x} \Fib(w, y)$ where $x > 0$."
%   "\item[\it Segment:] $[w, v]_x^y \coloneq \Fib(w, x) \cap \Fib(v, -y)$ where $x, y \geq 0$."
\begin{definition}[Task Relation]\label{def:task-relation}
A \textbf{task relation} on a nonempty set of world states $\worldstate$ over a temporal order $\D$ is any parameterised relation $w \taskto{x} u$ for $w, u \in \worldstate$ and $x \in \poscone$, extended to negative durations by the \emph{converse convention}
\[
  w \taskto{-x} u \quad \define \quad u \taskto{x} w
  \qquad \text{for } x \geq 0,
\]
determining the following for any world states $w, v \in \worldstate$ and durations $x, y \in D$:
\begin{itemize}
  \item[\emph{Fiber}:] $\Fib(w, x) \define \{\, u \in \worldstate : w \taskto{x} u \,\}$.
  \item[\emph{Cone}:] $\taskcone{w}{x} \define \bigcup_{|y| < x} \Fib(w, y)$ where $x > 0$.
  \item[\emph{Segment}:] $\taskseg{w}{v}{x}{y} \define \Fib(w, x) \cap \Fib(v, -y)$ where $x, y \geq 0$.
\end{itemize}
\end{definition}

\noindent
The fiber $\Fib(w, x)$ collects the world states reachable from $w$ by a task of duration exactly $x$; by the converse convention, negative-duration fibers run the relation backwards.
The cone $\taskcone{w}{x}$ collects the world states reachable forwards or backwards from $w$ within any duration of magnitude strictly less than $x$.
The segment $\taskseg{w}{v}{x}{y}$ collects the world states reachable forwards from $w$ in duration $x$ that also reach $v$ forwards in duration $y$: the candidate states lying between the two endpoints.
Fibers and segments are two separate classes of sets; in particular, a one-sided fiber does not count as a segment.

% Restates the JPL paper's \label{def:directed}. Verbatim source text:
%   "A nonempty family of sets $\mathcal{S}$ is \textit{directed} just in case
%    $S \subseteq S_1 \cap S_2$ for some $S \in \mathcal{S}$ whenever $S_1, S_2 \in \mathcal{S}$."
\begin{definition}[Directed Family]\label{def:directed}
A nonempty family of sets $\mathcal{S}$ is \textbf{directed} just in case $S \subseteq S_1 \cap S_2$ for some $S \in \mathcal{S}$ whenever $S_1, S_2 \in \mathcal{S}$.
\end{definition}

% Restates the JPL paper's \label{def:frame}: exactly FOUR axioms (Compositionality,
% Seriality, Limit, Saturation). Nullity is NOT an axiom -- it is the derived \ref{lem:nullity}
% below. Verbatim axiom text (sub-anchors def:frame#Compositionality, def:frame#Seriality,
% def:frame#Limit, def:frame#Saturation):
%   "\item[\it Compositionality:] $w \Rightarrow_{x + y} v$ if and only if $w \Rightarrow_x u$
%    and $u \Rightarrow_y v$ for some $u \in W$."
%   "\item[\it Seriality:] $w \Rightarrow_x u$ and $v \Rightarrow_x w$ for some $u, v \in W$."
%   "\item[\it Limit:] $\bigcap\limits_{x > 0} (w)_x = \set{w}$."
%   "\item[\it Saturation:] $\bigcap \mathcal{S} \neq \emptyset$ for any $\supseteq$-directed
%    family $\mathcal{S}$ of nonempty fibers and segments."
\begin{definition}[Task Frame]\label{def:frame}
A \textbf{task frame} is any structure $\taskframe = (\worldstate, \D, {\taskto{}})$ where $\worldstate$ is a nonempty set of world states, $\D = (D, +, 0, \leq)$ is a nontrivial totally ordered abelian group, and ${\taskto{}}$ is a task relation satisfying the following for $x, y \geq 0$:
\begin{itemize}
  \item[\emph{Compositionality}:] $w \taskto{x+y} v$ if and only if $w \taskto{x} u$ and $u \taskto{y} v$ for some $u \in \worldstate$.
  \item[\emph{Seriality}:] $w \taskto{x} u$ and $v \taskto{x} w$ for some $u, v \in \worldstate$.
  \item[\emph{Limit}:] $\bigcap_{x > 0} \taskcone{w}{x} = \{w\}$.
  \item[\emph{Saturation}:] $\bigcap \mathcal{S} \neq \emptyset$ for any $\supseteq$-directed family $\mathcal{S}$ of nonempty fibers and segments.
\end{itemize}
\end{definition}

\noindent
\emph{Compositionality} is a biconditional.
Read right to left, it asserts closure under composition: tasks of durations $x$ and $y$ performed in sequence compose to a single task of duration $x + y$.
Read left to right, it asserts interpolation: every task of duration $x + y$ passes through an intermediate world state that splits it into consecutive tasks of durations $x$ and $y$.
Both directions are load-bearing; in particular the interpolation direction is used directly in the proofs about constraint families of fibers and segments.
\emph{Seriality} provides, at every nonnegative duration, both a successor and a predecessor for every world state, so no world state is a temporal dead end in either direction.
\emph{Limit} excludes world states that are distinct from $w$ and yet reachable from $w$ within arbitrarily small durations; equivalently, $w$ is the only world state lying in every cone around $w$.
\emph{Saturation} is a directed-intersection condition: any $\supseteq$-directed family drawn from the nonempty fibers and segments has a common element.
It supplies the candidate world state needed to extend a partial history to one further time, and is consumed exactly at that step of the extension argument for possible worlds.

% Restates the JPL paper's derived \label{lem:nullity} (Nullity is a lemma, not an axiom).
% Verbatim source text:
%   "$w \Rightarrow_0 w$ for every world state $w \in W$ in every frame
%    $\F = \tuple{W, \D, \Rightarrow}$."
\begin{lemma}[Nullity]\label{lem:nullity}
$w \taskto{0} w$ for every world state $w \in \worldstate$ in every task frame $\taskframe = (\worldstate, \D, {\taskto{}})$.
\end{lemma}

\noindent
Nullity is \emph{derived} rather than postulated, and the derivation is choice-free: \emph{Seriality} at $x = 0$ supplies some $u$ with $w \taskto{0} u$; since $|0| < x$ for every $x > 0$, this $u$ lies in every cone around $w$, and \emph{Limit} then forces $u = w$.

\begin{remark}
Composition over negative durations is likewise derived rather than postulated: take converses, compose on the positive cone via \emph{Compositionality}, and take the converse of the result.
Mixed-sign composition is not so much prohibited as inexpressible at the primitive level, since primitive durations are nonnegative.
\end{remark}

\subsection{Convex Histories}

A convex history is a function from times to world states that respects the task relation over a convex temporal domain.
Convex histories represent possible paths through the space of world states; the ones whose domain is total are the \textbf{possible worlds}.

\begin{definition}[Convex Domain]
A domain $\domain : D \to \text{Prop}$ is \textbf{convex} if whenever $a, c \in \domain$ with $a \le c$, every time $b$ with $a \le b \le c$ is also in $\domain$.
More precisely, for all $a, b, c : D$, if $\domain(a)$ and $\domain(c)$ and $a \le b \le c$, then $\domain(b)$.
Convexity ensures the domain has no temporal gaps.
\end{definition}

\begin{definition}[Convex History]
A \textbf{convex history} in a task frame $\taskframe$ is a dependent function $\tau : (x : D) \to \domain(x) \to \worldstate$ where $\domain : D \to \text{Prop}$ is a convex subset of $D$ and $\tau(x) \taskto{y-x} \tau(y)$ for all times $x, y : D$ with $\domain(x)$, $\domain(y)$, and $x \le y$.
We write $\histories_{\taskframe}$ for the \textbf{possible worlds} over frame $\taskframe$: the convex histories whose domain is total, so that $\domain(x)$ holds for every $x : D$.
\end{definition}

\subsection{Task Models}

A task model extends a task frame with an interpretation function that assigns truth values to sentence letters at world states.
\textbf{Propositions} are subsets of $\worldstate$ representing instantaneous ways for the system to be.
Sentence letters express propositions, which can be realized by zero or more world states.
World states themselves are specific configurations of the total system at an instant.

\begin{definition}[Task Model]
A \textbf{task model} defined over a frame $\taskframe$ is a pair $\model = (\taskframe, I)$ where the \textbf{interpretation function} $I : \worldstate \to \text{String} \to \text{Prop}$ assigns to each world state $w : \worldstate$ and sentence letter $p : \text{String}$ a truth value $I(w, p) : \text{Prop}$.
We write $I(w, p)$ to indicate that sentence letter $p$ is true at $w$.
\end{definition}

\subsection{Truth Conditions}

Truth is evaluated relative to a model $\model$ providing the interpretation, a possible world $\tau$ representing a complete path through the space of world states, and a time $x : D$.
Whereas the model fixes the interpretation of the language, the contextual parameters $\tau$ and $x$ determine the truth value of every sentence of the langauge.

\begin{definition}[Truth]
For model $\model$, history $\tau : \histories_{\taskframe}$, and time $x : D$:
\begin{align*}
  \model, \tau, x \vDash p &\Iff x \in \domain(\tau) \text{ and } I(\tau(x), p) \\
  \model, \tau, x \nvDash \falsum \\
  \model, \tau, x \vDash \varphi \imp \psi &\Iff
    \model, \tau, x \nvDash \varphi \text{ or } \model, \tau, x \vDash \psi \\
  \model, \tau, x \vDash \nec\varphi &\Iff
    \model, \sigma, x \vDash \varphi \text{ for all } \sigma : \histories_{\taskframe} \\
  \model, \tau, x \vDash \allpast\varphi &\Iff
    \model, \tau, y \vDash \varphi \text{ for all } y : D \text{ where } y < x \\
  \model, \tau, x \vDash \allfuture\varphi &\Iff
    \model, \tau, y \vDash \varphi \text{ for all } y : D \text{ where } x < y
\end{align*}
\end{definition}

\noindent
The modal operator $\nec$ quantifies over all possible worlds $\sigma : \histories_{\taskframe}$ at the current time $x : D$.
The temporal operators $\allpast$ and $\allfuture$ quantify over all earlier and later times $y : D$.


\subsection{Time-Shift}

The time-shift operation is used to establish the \textbf{perpetuity principles}:
\begin{itemize}
  \item P1: $\nec\varphi \imp \always\varphi$ (what is necessary is always the case)
  \item P2: $\sometimes\varphi \imp \poss\varphi$ (what is sometimes the case is possible)
\end{itemize}
It is natural to assume that whatever is necessary is always the case, or equivalently, whatever is sometimes the case is possible.
Time-shift enables proofs of the bimodal axioms MF ($\nec\varphi \imp \nec\allfuture\varphi$) and TF ($\nec\varphi \imp \allfuture\nec\varphi$) which together imply the perpetuity principles.

\begin{definition}[Time-Shift]
For $\tau, \sigma \in \histories_{\taskframe}$ and $x, y : D$, possible worlds $\tau$ and $\sigma$ are \textbf{time-shifted from $y$ to $x$}, written $\tau \approx_y^x \sigma$, if and only if there exists an order automorphism $\bar{a} : D \to D$ where $y = \bar{a}(x)$, $\domain_\sigma = \bar{a}^{-1}(\domain_\tau)$, and $\sigma(z) = \tau(\bar{a}(z))$ for all $z \in \domain_\sigma$.
\end{definition}

\noindent
Time-shifting preserves the essential structure of histories:

\begin{theorem}[Convexity Preservation]
If $\tau$ has a convex domain and $\tau \approx_y^x \sigma$, then $\sigma$ has a convex domain.
\end{theorem}

\begin{theorem}[Task Preservation]
If $\tau$ respects the task relation and $\tau \approx_y^x \sigma$, then $\sigma$ respects the task relation.
\end{theorem}

\subsection{Logical Consequence and Validity}

Logical consequence and validity are defined uniformly across all temporal types, frames, models, histories, and times.

\begin{definition}[Logical Consequence]
A formula $\varphi$ is a \textbf{logical consequence} of $\Gamma$ (written $\Gamma \models \varphi$) just in case for every temporal type $D : \text{Type}$, frame $\taskframe : \text{TaskFrame}(D)$, model $\model : \text{TaskModel}(\taskframe)$, history $\tau \in \histories_{\taskframe}$, and time $x : D$, if $\model, \tau, x \models \psi$ for all $\psi \in \Gamma$, then $\model, \tau, x \models \varphi$.
\end{definition}

\begin{definition}[Validity]
A formula $\varphi$ is \textbf{valid} (written $\models \varphi$) just in case $\varphi$ is a logical consequence of the empty set: $\varnothing \models \varphi$.
Equivalently, $\varphi$ is true at every model-history-time triple.
\end{definition}

\begin{definition}[Satisfiability]
A context $\context$ is \textbf{satisfiable} in temporal type $D : \text{Type}$ if there exist a frame $\taskframe : \text{TaskFrame}(D)$, model $\model : \text{TaskModel}(\taskframe)$, history $\tau \in \histories_{\taskframe}$, and time $x : D$ such that $\model, \tau, x \models \psi$ for all $\psi \in \context$.
\end{definition}

\begin{theorem}[Monotonicity]
If $\context \subseteq \Delta$ and $\context \models \varphi$, then $\Delta \models \varphi$.
\end{theorem}

\end{document}
